\section{Introdução}\label{sec:introducao}

% ─── Beat 1: O problema. ─────────────────────────────────────────────
% Jogos em nuvem: cliente fino, servidor caro, rede no meio.
% Requisitos duros: RTT < 100 ms, perda < 1\%.
% A rede doméstica não os respeita o tempo todo.

% ─── Beat 2: Por que controle adaptativo de taxa de bits importa. ────
% Sem ABR: estouro de buffer (bufferbloat), travamentos, recuperação
% via I-frame (IDR) — punição visual severa.
% O servidor é o lugar natural para o controle: ele controla o codec.

% ─── Beat 3: O Sunshine e o problema das plataformas comerciais. ────
% Plataformas comerciais (GeForce NOW, Xbox Cloud Gaming, PlayStation
% Plus) resolvem o problema às custas de lock-in: assinatura mensal,
% catálogo restrito, hardware proprietário no data center, sem controle
% do usuário sobre qualidade/latência. Alternativa open-source: o
% \emph{Sunshine}, servidor de streaming compatível com o protocolo
% NVIDIA GameStream (cliente Moonlight), permite que o usuário rode o
% próprio jogo na própria máquina e estreame para qualquer dispositivo
% leve em sua rede ou via internet. Hoje, opera com taxa de bits fixa
% negociada no início da sessão — sem reação a degradação de rede em
% curso, deixando a experiência refém da estabilidade do link.

% ─── Beat 4: A contribuição. ────────────────────────────────────────
% Implementamos um controlador AIMD com primitivas do CUBIC
% (Ha, Rhee, Xu, 2008) — curva cúbica de recuperação e convergência
% rápida — apenas com sinais observáveis no servidor (RTT, sem
% cooperação do cliente). Avaliamos contra (a) base sem ABR e
% (b) AIMD de duas fases sob seis cenários de degradação reproduzíveis.

% ─── Beat 5: Resultado principal. ───────────────────────────────────
% Redução de até 43\% no pico de RTT e até 60\% no tempo total de
% travamento, ao custo de um leve aumento na latência mediana — uma
% troca alinhada às prioridades de jogabilidade discutidas em
% \cite{alos2019,abragame}.

% ─── Beat 6: Estrutura do trabalho. ─────────────────────────────────
\subsection{Estrutura do Trabalho}
% Cap. \ref{sec:referencial} apresenta o referencial teórico e os
% trabalhos correlatos. Cap. \ref{sec:contribuicao} descreve o
% controlador proposto. Cap. \ref{sec:metodologia} detalha o testbed
% e o protocolo experimental. Cap. \ref{sec:resultados} apresenta os
% resultados. Cap. \ref{sec:conclusao} conclui e aponta trabalhos
% futuros.
